%_____Zusammenfassung, Ausblick_________________________________
\chapter{Conclusion}
\label{chap:conclusion}

This thesis analyzed the combination of Lie group
variational integrators (LGVIs) with sparse Moment-SOS semidefinite relaxations
for multi rigid-body trajectory optimization. 

Firstly, Chapter \ref{sec:introduction} 
formulated the general optimization problem and identified the research gap of 
semidefinite relaxation providing global optimal solutions
for some low-dimensional rigid-body systems, however, not analyzed for 
different complex multibody systems.
Thereby, accuracy of lower bounds, relaxation tightness and feasible extraction become
problem dependent as the dynamical complexity increases. This question was
examined for a single rigid body and a constrained, underactuated multibody
system.

In order to derive LGVIs, 
Chapter \ref{chapter:preliminaries_lie_groups} briefly introduced the matrix 
Lie groups, Lie algebras, and local
group variations required for the following derivations.
%%%%%%

Chapter \ref{chapter:continous_mechanics_of_rigid_bodies_on_lie_groups} 
established the continuous mechanical foundation. The
Euler-Lagrange equations on a matrix Lie group were derived from Hamilton's
principle using left-trivialized variations \eqref{eq:continuous_euler_lagrange_lie_group}.
Additionally, the Legendre transformation was
used to express the dynamics in Hamiltonian form \eqref{eq:continuous_hamilton_equations_lie_group}. 
This framework was applied to
the unforced 3D pendulum on $\mathrm{SO}(3)$, for which the continuous equations
were derived \eqref{eq:continuous_3D_pendulum_euler_lagrange}. 
To validate the LGVI properties of the 3D pendulum,
the momentum conservation associated with rotations about gravity was
verified \eqref{eq:continuous_3d_pendulum_validation} using Noether's theorem.

Afterwards, Chapter \ref{chapter:discrete_mechanics_for_rigid_bodies_on_lie_groups} 
transferred the continous lagrangian variational construction to discrete time. 
Instead of
discretizing the continuous equations of motion directly, the action integral was
replaced by a discrete action sum \eqref{eq:discrete_action_sum}. 
This yields to the discrete Euler-Lagrange \eqref{eq:discrete_euler_lagrange_lie_group} 
and Hamiltonian equations \eqref{eq:discrete_hamilton_equations_lie_group}, 
discrete Legendre transformations \eqref{eq:discrete_legendre_transformations}, 
and their forced extension \eqref{eq:discrete_forced_Euler_Lagrange}. 
A computational approach was introduced to solve the implicit discrete equations of motion. 
Thereby, multiplicative kinematic updates preserve the Lie-group structure, 
while a Cayley parametrization and Newton iteration solved the implicit rotation 
update. In corespondance, to the continous case, the discrete 
unforced and controlled 3D-pendulum LGVI were derived \eqref{eq:discrete_3D_pendulum_EL}. 
This single rigid-body formulation was
extended to a multi-body framework for deriving LGVI for coupled rigid bodies. 
Therefore, the Acrobot LGVI using maximal coordinates was derived as an example 
\eqref{eq:acrobot_full_del_shifted}. Its two center-of-mass positions and
$\mathrm{SO}(2)$ attitudes were coupled by holonomic constraints and
Lagrangian multipliers.

To optimize these models trajectories, Chapter \ref{chapter:certifiable_trajectory_optimization} 
transforms these discrete dynamical models into finite-horizon polynomial
optimization problems (POPs). However, as these equations are nonconvex
semidefinite programming is introduced to convert these into convex optimization problems.
Thereby the Moment-SOS hierarchy and sparse relaxations are introduced. 
Afterwards, the POP of the 3D-pendulum is defined on
$\mathrm{SO}(3)$ \eqref{eq:controlled_3d_pendulum_pop}. 
For the Acrobot, the derived LGVI was reduced to a smaller variable size optimization problem.
Thereby, the holonomic constraints eliminated
the positions and the velocity variables, which were expressed through rotations, 
reducing the decision vector from $17N+3$ to $13N-1$ variables. Still, 
the quadratic structure of the dynamics was preserved. The resulting POP was defined in \eqref{eq:controlled_acrobot_pop}.
In addition, chain-like cliques \eqref{eq:acrobot_full_self_interior_cliques} 
were compared with a multibody-specific
clique construction \eqref{eq:acrobot_hybrid_self_cliques} 
to seperate moment matrices for computational efficiency.
Finally, the reduced Acrobot POP was embedded in an SDP-based
model predictive control (MPC) loop that applied only the first extracted control
after each SDP optimization itteratively.

In Consequence, Chapter \ref{chapter:evaluation} 
evaluated the resulting models and optimization problems. 
Firstly, a numerical simulation of the 3D-pendulum LGVI 
and the Acrobot LGVI was implemented to verify the derived equations of motion.
As result, the unforced 3D-pendulum simulation remained on $\mathrm{SO}(3)$, 
preserved the momentum, and conserved the energy over a long time horizon. 
Similarly, the Acrobot simulations preserved both $\mathrm{SO}(2)$ attitudes and 
the base and elbow constraints. 
Secondly, the POPs were solved using SDP relaxations and evaluated for 
feasibility and optimality. All five evaluated 3D-pendulum relaxations 
were numerically rank one, gave feasible extractions, and were certified for 
their discrete POPs (see Table \ref{tab:pendulum_3d_sdp_statistics}). 
For the 3D-pendulum minimum degree chordal extension
was used for the clique construction. Moreover, the Acrobot $0^\circ$ and $45^\circ$ targets were also 
certified using chain-like cliques and multibody-specific clique constructions. 
In contrast, the $90^\circ$,
$135^\circ$, and $180^\circ$ semidefinite relaxation cases produced higher-rank moment
matrices and infeasible extractions (see Table \ref{tab:acrobot_sdp_statistics_full}-\ref{tab:acrobot_sdp_statistics_hybrid}). 
Therefore, no global certificate was obtained for these cases. Nevertheless,
the closed loop MPC implementation with repeated
first-control extraction generated feasible simulated Acrobot swing-ups. With
the $15\,\mathrm{N\,m}$ bound, the two matched time-scale cases each stabilized
$94$ of $100$ sampled states, compared with $34$ of $100$ for the mismatched
case (see Fig.\ref{fig:acrobot_basin_comparison}). 

Finally, Chapter \ref{chap:discussion} interpreted these findings in relation to the research question and
the state of the art literature. Additional, possible applications of optimization and learning-based methods were identified.

In connection to the research question the results provide following analysis. 
The LGVI
dynamics of both considered systems can be represented by low-degree polynomial
constraints while retaining their Lie-group structure. The resulting POPs can be solved by sparse Moment-SOS relaxations, 
which can provide certified global optima for some low-dimensional
rigid-body problems. Here, Sparse second-order Moment-SOS relaxation recovered feasible trajectories that
were certified as globally optimal for the tight 3D-pendulum problems and the two
simpler Acrobot optimization problems. However, the Acrobot cases of 
$\theta_{des} \geq 90^\circ$ prevent this result from being
generalized to complex multibody motions. 
Also, sparsity reduces the SDP computation time, 
but does not guarantee tightness or feasible extraction.
The approach of an closed-loop SDP-MPC produces feasible trajectories 
for even complex optimization problems, for example the stabilization of 
an underactuated Acrobot,
when the prediction was not rank one. However, it did not globally certify 
these solutions.

\paragraph{Limitations}

This thesis does not consider contact, collision avoidance and obstacles. Adittionally,
hardware experiments were not performed.

Computational scaling is the strongest practical limitation because the moment
matrix dimension grows with the variables in each clique. Without
sparsity, the reported problems would require matrices of size
$91{,}378\times91{,}378$ and $304{,}590\times304{,}590$. Even with sparsity,
MOSEK required about $56$ to $69$ min for one 3D-pendulum problem,
$15$--$31$ min for the special-multibody Acrobot cliques, and $45$--$113$ min
for the complete chain-like cliques on eight CPU cores. The successful initial
$(80^\circ,80^\circ)$ Acrobot MPC computations required $3.1$, $27.3$, and $17.8$ hours
in Cases I--III. The implementation is therefore an offline study. In general, larger
time horizons, systems, or coupled constraints are increasing the computational
expense significantly.

\paragraph{Future Work}

A first direction is to reduce the polynomial representation before constructing
the SDP. For spatial rotations, a unit-quaternion formulation should be compared
with the rotation-matrix formulation used here. The relevant question is not only
whether fewer scalar attitude variables are obtained, but whether the complete
quaternion POP produces smaller sparse moment blocks while preserving suitable
polynomial dynamics and reliable extraction. This comparison should include the
resulting polynomial degree, constraints, clique structure, numerical conditioning,
and relaxation tightness.

The framework should next be transferred to additional multibody benchmarks. A
cart--pole system provides a useful intermediate problem containing translational
and rotational motion. Its LGVI and reduced POP could be derived, its temporal
and problem-specific sparsity identified, and the generic and user-defined clique
constructions compared. The same SDP--MPC procedure could then be applied to
determine whether the first-control strategy remains effective. This would provide
a controlled progression toward systems with more links, followed by obstacle,
contact, uncertainty, and hardware experiments already identified as open in this
thesis.

The computational backend should also be reconsidered. A first-order GPU solver
such as CuADMM could replace the interior-point solver MOSEK and substantially
reduce solution time, although potentially lower numerical accuracy
would have to be evaluated carefully whenever a certificate is claimed. The solver
itself would not make a fixed relaxation tighter. However, faster solution of the
SDP could make larger cliques or higher relaxation orders affordable, which may
yield tighter bounds for the difficult Acrobot cases, and could reduce the cost of
longer horizons. Warm starting consecutive MPC problems and reusing their
unchanged polynomial and sparsity structure should be investigated for the same
reason.

Finally, future work could examine adaptive and iterative extraction. Rank metrics
could identify locally tight moment blocks, after which only variables supported by
verified rank-one blocks would be fixed and the remaining relaxation solved again.
Clique merging or a higher relaxation order could be applied selectively where the
rank condition fails. Such a procedure may help construct a feasible full-horizon
trajectory, but fixing local decisions alone is not a certificate of global optimality.
A global claim would still require satisfaction of all constraints of the original POP
and agreement between the objective of the completed feasible trajectory and a
valid lower bound for that original problem. In parallel, hard terminal constraints
or terminal sets should be evaluated so that certification refers to reaching the
target rather than only minimizing a soft tracking cost. For the closed-loop method,
a systematic study of matched prediction and execution time scales, together with
terminal stability conditions, would be required before extending the approach
from empirical stabilization to a certified controller.


\begin{comment}
	
Do not just leave it at the discussion: discuss what you/the reader can learn from the results. Draw some real conclusions. Separate discussion/interpretation of the results clearly from the conclusions you draw from them. (So-called ``conclusion creep'' tends to upset reviewers. It means surrendering your scientific objectivity.) Identify all shortcomings/limitations of your work, and discuss how they could be fixed (``future work''). It is not a sign of weakness of your work if you clearly analyze and state the limitations. Informed readers will notice them anyway and draw their own conclusions if not addressed properly.

\vspace{\baselineskip}
Recap: do not stick to this structure at all cost. Also, remember that the thesis must be:

\begin{itemize}
	\item honest, stating clearly all limitations
	\item self--contained, do not write just for the locals, do not assume that the reader has read the same literature as you, do not let the reader work out the details for themselves
\end{itemize}

This chapter is followed by the list of figures and the bibliography. If you are using acronyms, listing them (with the expanded full name) before the bibliography is also a good idea. The acronyms package helps with consistency and an automatic listing.
\end{comment}
